

\sfss{Context}

    \begin{describe}[black]
        Describe the process or a system at a high level. (Can reference a process card for more details. Can call out what is unspecified.)
    \end{describe}
    \dscplaceholder

    \begin{describe}[black]
        What are other systems that this process depends on or interacts with, which impact its success? (e.g., sortition data, or user or citizen authentication systems)
    \end{describe}
    \dscplaceholder
    
\sfss{Process Quality}
\sfsss{Representation}

    \begin{describe}[black]
        The extent to which key decisions are representative of the constituent population.
    \end{describe}
    \dscplaceholder

    \begin{evaluate}[black]
        To what extent: (1) is there sufficient representation at critical parts of the process, including (a) proposing decisions, and (b) making ultimate decisions? (2) are there barriers leading to bias in representation?
    \end{evaluate}
    \dscplaceholder
    
\sfsss{Informedness}

    \begin{describe}[black]
        The extent to which critical information is taken into account for decision-making.
    \end{describe}
    \dscplaceholder
    
    \begin{evaluate}[black]
        To what extent: (1) is critical context incorporated into decision-making about tradeoffs and consequences of different decisions? (2) is this sourced from (a) domain expertise, (b) the existing authorities, who may have extensive context, (c) a broad diversity of constituents, (d) the most impacted stakeholders, and (e) the powerful stakeholders, whose incentives are critical to having the decision ``stick''?
    \end{evaluate}
    \dscplaceholder
    
\sfsss{Deliberation}
    
    \begin{describe}[black]
        The extent to which decisions are considered and deliberative (rather than superficial and reactive).
    \end{describe}
    \dscplaceholder
    
    \begin{evaluate}[black]
        To what extent are those involved: (1) able to (and supported to) move from shallower to deeper goals and values? (2) able to (and supported to) collaborate where necessary? (3) able to address issues within the available time?
    \end{evaluate}
    
\sfsss{Substantiveness}
    
    \begin{describe}[black]
        The extent to which decisions are substantive (e.g., actionable, consequential) rather than nonsubstantive (e.g., vague, simplistic, inconsequential).
    \end{describe}
    \dscplaceholder
    
    \begin{evaluate}[black]
        To what extent: (1) is the decision directly actionable and implementable? (2) does the decision meaningfully address the issues? (3) does the decision grapple with the necessary levels of complexity? (4) is uncertainty appropriately managed and accounted for? (5) are risks to implementability accounted for?
    \end{evaluate}
    
\sfsss{Robustness}
    
    \begin{describe}[black]
        The extent to which the process is robust to suboptimal conditions or adversarial or strategic behavior.
    \end{describe}
    \dscplaceholder
    
    \begin{evaluate}[black]
        To what extent is the process or system vulnerable to: (1) suboptimal conditions or broken assumptions? (e.g., low turnout, larger power asymmetries) (2) strategic behavior and manipulation? (3) false claims? (e.g., of manipulation)
    \end{evaluate}
    \dscplaceholder
    
\sfsss{Legibility}
    
    \begin{describe}[black]
        The extent to which the processes and decisions are accessible, understandable, and verifiable.
    \end{describe}
    \dscplaceholder
    
    \begin{evaluate}[black]
        To what extent is information (a) accessible, (b) understandable, (c) verifiable about the: (1) processes/systems used to make decisions? (2) the execution of these processes? (3) decisions being made (4) reasons and inputs feeding into decisions?
    \end{evaluate}
    \dscplaceholder
    
\sfss{Delegation}
\sfsss{Integration}
    
    \begin{describe}[black]
        The extent to which the authority integrates the democratic process into its operations.
    \end{describe}
    \dscplaceholder
    
    \begin{evaluate}[black]
        To what extent is the authority structuring its internal communications and operations to effectively: (1) provide critical context to the democratic process / system? (2) integrate democratic process outputs in its actions? (3) trigger democratic processes when/if required?
    \end{evaluate}
    \dscplaceholder
    
\sfsss{Ability to bind}
    
    \begin{describe}[black]
        The extent to which the authority is able to technically and legally bind itself to democratic decisions.
    \end{describe}
    \dscplaceholder
    
    \begin{evaluate}[black]
        To what extent can the unilateral authority bind itself to acting in accordance with the democratic decision: (1) technically? (2) legally? (e.g., has developed the needed technical and/or legal infrastructure for binding)
    \end{evaluate}
    \dscplaceholder
    
\sfsss{Commitment}
    
    \begin{describe}[black]
        The extent to which the unilateral authority commits to acting in accordance with the democratic decision.
    \end{describe}
    \dscplaceholder
    
    \begin{evaluate}[black]
        To what extent has the unilateral authority committed to acting in accordance with the democratic decision: (1) internally? (2) privately? (3) publicly?  (regardless of their ability to bind)
    \end{evaluate}
    \dscplaceholder
    
\sfss{Trust}
\sfsss{Awareness}
    
    \begin{describe}[black]
        The extent to which the relevant public is aware of the democratic process.
    \end{describe}
    \dscplaceholder
    
    \begin{evaluate}[black]
        To what extent is the relevant public aware: (1) that the democratic system exists? (2) how it works? (3) what it is being used for? (4) how they can be involved?
    \end{evaluate}
    \dscplaceholder

\sfsss{Participation}
    
    \begin{describe}[black]
        The extent to which the relevant public is willing to participate in the process.
    \end{describe}
    \dscplaceholder
    
    \begin{evaluate}[black]
        To what extent is the relevant public: (1) willing to participate? (2) able to participate? (3) appropriately compensated for participating? (4) actually participating?
    \end{evaluate}
    \dscplaceholder
    
\sfsss{Accountability}
    
    \begin{describe}[black]
        The extent to which there are external watchdogs and accountability structures monitoring the execution of the democratic process and the implementation of its outputs.
    \end{describe}
    \dscplaceholder
    
    \begin{evaluate}[black]
        To what extent are: (1) there well-understood lines of oversight and accountability? (2) sufficiently influential/powerful organizations focused on holding authorities to their promised levels of democratic involvement? (3) authorities and democratic systems responsive to such accountability mechanisms?
    \end{evaluate}
    \dscplaceholder
    
\sfsss{Buy-in}
    
    \begin{describe}[black]
        The extent to which the relevant public and key stakeholders buy into the process and its legitimacy.
    \end{describe}
    \dscplaceholder
    
    \begin{evaluate}[black]
        To what extent are the relevant public and key stakeholders accepting of the legitimacy of: (1) the system/process? (2) of the decision?
    \end{evaluate}
    \dscplaceholder

\clearpage
